\documentclass[../../../../main.tex]{subfiles}
\begin{document}

直円錐形のグラスに水が満ちている．
水面の円の半径は1，深さも1である．

\begin{figure}[htbp]
  \centering
  \begin{tikzpicture}[scale=1.5, >=stealth]
    \coordinate (V) at (0, -1.2);
    \coordinate (O) at (-1.5, 0);
    \coordinate (A) at (1.2, 1.8);

    \coordinate (W1) at (-1.5, 0);
    \coordinate (W2) at (0.8, 0);

    \fill[black!15] (V) -- (W1) -- (W2) -- cycle;

    \draw[thick] (W1) -- (W2);
    \draw[thick] (A) -- (V) -- (O);
    \draw[thick] (O) -- (A);

    \draw (O) +(0:0.5) arc (0:30:0.5);
    \node at ($(O)+(15:0.7)$) {$\alpha$};
  \end{tikzpicture}
\end{figure}

\begin{enumerate}
  \item このグラスを右の図のように角度 $\alpha$ だけ傾けたとき，
  できる水面は楕円である．
  この楕円の中心からグラスのふちを含む平面までの距離 $\ell$ と，
  楕円の長半径 $a$ および短半径 $b$ を，
  $m=\tan\alpha$ で表せ．
  ただし楕円の長半径，短半径とは，それぞれ長軸，短軸の長さの $\displaystyle\frac{1}{2}$ のことである．

  \item 傾けたときこぼれた水の量が，
  最初の水の量の $\displaystyle\frac{1}{2}$ であるとき，
  $m=\tan\alpha$ の値を求めよ．
  ただしグラスの円錐の頂点から，
  新しい水面までの距離を $h$ とするとき，
  残った水の量は，
  $\displaystyle\frac{1}{3}\pi abh$ に等しいことを用いよ．
\end{enumerate}

\end{document}
